
%%
%% Part 6 - Flow Matching Implementation and Evaluation
%%
%% Building on the theoretical foundations laid in Part 5, this chapter
%% describes how to implement Flow Matching for one-dimensional time
%% series and evaluates its performance on synthetic financial data.
%% We introduce the transformer-based velocity network, the training
%% objective, sampling via numerical ODE solvers and a suite of
%% evaluation metrics.

\section{Flow Matching in Practice}

This chapter translates the Flow Matching theory into a working
algorithm for generating financial time series.  We present our
network architecture, describe the training procedure, outline the
sampling routine and analyse the resulting synthetic data.

\subsection{Network Architecture}

We model the velocity field $v_\theta(x,t)$ using a transformer
encoder that operates on sequences of returns.  Each input is a
matrix of shape $[\text{batch}, L, F]$, where $L$ is the sequence
length and $F$ the number of features (here $F=1$ for univariate
returns).  A sinusoidal time embedding of $t$ is added to every
sequence position, and then multiple layers of multi-head
self-attention and feed-forward networks are applied.  The output is
projected back to the original dimensionality to yield a velocity
vector for each time step in the sequence.

\begin{lstlisting}[language=Python, caption=Time Series Flow Matching Model]
class TimeSeriesFlowMatching(nn.Module):
    def __init__(self, input_dim=1, hidden_dim=128, num_layers=4, num_heads=8):
        super().__init__()
        self.time_embed = nn.Sequential(
            SinusoidalEmbedding(hidden_dim),
            nn.Linear(hidden_dim, hidden_dim),
            nn.GELU(),
            nn.Linear(hidden_dim, hidden_dim)
        )
        self.input_proj = nn.Linear(input_dim, hidden_dim)
        self.layers = nn.ModuleList([
            nn.TransformerEncoderLayer(
                d_model=hidden_dim,
                nhead=num_heads,
                dim_feedforward=4*hidden_dim,
                dropout=0.1,
                batch_first=True
            ) for _ in range(num_layers)
        ])
        self.output_proj = nn.Linear(hidden_dim, input_dim)
        self.norm = nn.LayerNorm(hidden_dim)

    def forward(self, x, t):
        batch_size, seq_len, _ = x.shape
        t_embed = self.time_embed(t.unsqueeze(-1))
        t_embed = t_embed.unsqueeze(1).expand(-1, seq_len, -1)
        h = self.input_proj(x) + t_embed
        for layer in self.layers:
            h = layer(h)
        velocity = self.output_proj(self.norm(h))
        return velocity
\end{lstlisting}

\subsection{Conditional Flow Matching Loss}

To train the model, we sample pairs $(x_0,x_1)$ of noise and data
sequences and a random time $t\sim\mathrm{Uniform}(0,1)$.  We
interpolate between the two sequences via $x_t = (1-t)x_0 + tx_1$ and
compute the target velocity $v_{\text{target}} = x_1 - x_0$.  The
loss is the mean squared error between the predicted and target
velocity:
\begin{lstlisting}[language=Python, caption=CFM Loss]
def compute_cfm_loss(model, x0, x1, t):
    t_view = t.view(-1, 1, 1)
    xt = (1.0 - t_view) * x0 + t_view * x1
    v_target = x1 - x0
    v_pred = model(xt, t)
    loss = F.mse_loss(v_pred, v_target)
    return loss
\end{lstlisting}

Training proceeds over epochs, sampling new pairs from the dataset and
optimising with Adam.  Typical hyperparameters are listed in
Table~\ref{tab:fmparams}.

\begin{table}[H]
  \centering
  \begin{tabular}{lc}
    \toprule
    Parameter & Value \\
    \midrule
    Sequence Length & 100 \\
    Hidden Dimension & 128 \\
    Number of Layers & 4 \\
    Number of Heads & 8 \\
    Batch Size & 32 \\
    Learning Rate & $10^{-4}$ \\
    Training Epochs & 100 \\
    ODE Steps & 100 \\
    \bottomrule
  \end{tabular}
  \caption{Flow Matching training parameters.  Data are normalised
    globally to mean zero and unit variance.}
  \label{tab:fmparams}
\end{table}

\subsection{Sampling via ODE Integration}

After training, we generate new sequences by solving the ODE
\(dx/dt = v_\theta(x,t)\) from $t=0$ to $t=1$.  We initialise $x(0)$
with Gaussian noise and approximate the solution using Euler's method
or higher-order solvers.  For a simple Euler scheme with $N$ steps:
\begin{lstlisting}[language=Python, caption=Euler Sampling for Flow]
def sample_flow(model, x0, num_steps=100):
    dt = 1.0 / num_steps
    xt = x0.clone()
    for i in range(num_steps):
        t = i * dt
        t_tensor = torch.full((x0.shape[0],), t, device=x0.device)
        vt = model(xt, t_tensor)
        xt = xt + vt * dt
    return xt
\end{lstlisting}

Using more sophisticated solvers such as Heun's method or Runge–Kutta
improves accuracy at the expense of additional evaluations of
$v_\theta$.  We set $N=100$ in our experiments, which balances
quality and computational cost.

\subsection{Evaluation Metrics}

We evaluate Flow Matching using several metrics:
\begin{itemize}
  \item \textbf{Training and validation loss.}  We monitor the mean
    squared error on a held-out test set to detect overfitting.
  \item \textbf{Sliced Wasserstein Distance (SWD).}  By projecting
    both real and generated sequences onto random directions and
    computing one-dimensional Wasserstein distances, we approximate
    the Wasserstein distance in high dimensions.
  \item \textbf{Autocorrelation Function (ACF).}  We compute the ACF of
    real and generated returns for lags up to 20 and measure the mean
    squared or absolute difference.
  \item \textbf{Statistical Moments.}  We compare mean, variance,
    volatility (standard deviation), skewness and kurtosis of the real
    and synthetic returns.
\end{itemize}

\subsection{Empirical Results}

\paragraph{Loss curves.}  Figure~\ref{fig:fm_losses} shows the
training and test loss over 50 epochs.  The losses decrease smoothly,
with the test loss closely tracking the training loss, indicating
little overfitting.

\begin{figure}[H]
  \centering
  \fbox{\parbox{0.8\linewidth}{\centering
    \textit{Plot Placeholder:}\newline
    A line chart depicting training loss (blue) and test loss (orange)
    over epochs.  Both curves decrease to approximately 1.5 by epoch
    50.}}
  \caption{Training and validation loss for Flow Matching.}
  \label{fig:fm_losses}
\end{figure}

\paragraph{Synthetic sequences.}  Figure~\ref{fig:fm_synth} displays
several generated sequences alongside real data.  The synthetic
returns exhibit similar amplitudes, roughness and volatility
clustering to the ground truth.

\begin{figure}[H]
  \centering
  \fbox{\parbox{0.9\linewidth}{\centering
    \textit{Plot Placeholder:}\newline
    A set of line plots comparing real log-returns (blue) and
    generated log-returns (green dashed).  The patterns and
    variability are visually similar.}}
  \caption{Comparison of real and Flow Matching generated sequences.}
  \label{fig:fm_synth}
\end{figure}

\paragraph{Distribution comparison.}  Histogram and kernel density
estimations (Figure~\ref{fig:fm_dist}) reveal that the generated
returns closely match the central part of the real distribution but
underestimate extreme tails.  A Q-Q plot confirms that heavy-tailed
behaviour is not fully captured.

\begin{figure}[H]
  \centering
  \fbox{\parbox{0.9\linewidth}{\centering
    \textit{Plot Placeholder:}\newline
    Panels showing histograms and KDE curves of real and generated
    returns, a Q-Q plot and CDF comparison.  The central mass aligns
    well while the tails differ.}}
  \caption{Distribution comparison between real and generated data.}
  \label{fig:fm_dist}
\end{figure}

\paragraph{Statistical metrics.}  Table~\ref{tab:stats} lists basic
statistics of the real and synthetic data.  Mean and variance are
accurately reproduced, volatility matches within 1\% and skewness is
close to zero.  However, kurtosis is underestimated, reflecting
lighter tails.

\begin{table}[H]
  \centering
  \begin{tabular}{lcccc}
    \toprule
    Statistic & Real & Generated & Absolute Error & Relative Error \\
    \midrule
    Mean & 0.003770 & 0.000188 & 0.003582 & — \\
    Variance & 0.927920 & 0.909990 & 0.017930 & 1.9\% \\
    Volatility & 0.963286 & 0.953934 & 0.009352 & 0.97\% \\
    Skewness & 0.3667 & 0.0344 & 0.3323 & — \\
    Kurtosis & 2.2517 & 0.0784 & 2.1733 & — \\
    \bottomrule
  \end{tabular}
  \caption{Comparison of statistical moments for real and Flow
    Matching generated data.  The kurtosis discrepancy indicates
    underestimation of tail heaviness.}
  \label{tab:stats}
\end{table}

\paragraph{SWD and ACF.}  We compute a Sliced Wasserstein Distance of
approximately 0.0985 between the real and synthetic distributions,
indicating a good overall match.  The autocorrelation function of the
returns is near zero for both real and generated data, with an ACF
MSE of 0.0033.

\subsection{Discussion}

Flow Matching successfully captures the first two moments and the
volatility structure of the data with a simple linear path and a
transformer-based velocity field.  The deterministic nature of the
sampling makes generation efficient.  The main limitation is the
underestimation of extreme events, as manifested in the low kurtosis.
Future work could explore non-linear paths or explicit tail
modelling to address this issue.

\subsection{Conclusion of Part 6}

In this chapter we have implemented Flow Matching for financial time
series, described the transformer velocity network, the training and
sampling procedures, and evaluated the synthetic data.  Flow
Matching offers an appealing alternative to stochastic diffusion
models, with faster sampling and comparable performance on key
statistics.  The next chapter synthesises the lessons learned and
presents a comprehensive evaluation and conclusion.
